\begin{explanationbox}
	\textbf{\textcolor{CExplanation}{一句话直觉}}
	两条直线垂直等价于它们法向量的数量积为零，即两条直线对应系数的乘积之和等于零。

	\medskip
	\textbf{\textcolor{CExplanation}{核心拆解}}
	\begin{description}[style=nextline, leftmargin=2.8em, labelsep=0.5em, itemsep=0.45em]
		\item[\textbf{要点一（法向量引入）：}]
			直线$l_1:A_1x+B_1y+C_1=0$的一个法向量是$(A_1,B_1)$，类似地$l_2$的法向量为$(A_2,B_2)$。
		\item[\textbf{要点二（垂直对应点积零）：}]
			两直线垂直等价于两法向量垂直；两法向量垂直$\Leftrightarrow$数量积$A_1A_2+B_1B_2=0$。
	\end{description}


	\medskip
	\textbf{\textcolor{CExplanation}{为什么法向量是 $(A,B)$？（两点代入，相减得到）}}
	设直线 $l:Ax+By+C=0$（$A,B$ 不同时为零），在直线上任取两个不同的点
	$P(x_1,y_1)$、$Q(x_2,y_2)$。因为两点都在直线上，所以
	\[
	\begin{aligned}
	Ax_1+By_1+C&=0,\\
	Ax_2+By_2+C&=0.
	\end{aligned}
	\]
	\textbf{把两个方程相减}，消去常数项 $C$：
	\[
	A(x_2-x_1)+B(y_2-y_1)=0.
	\]
	而直线的一个方向向量正是 $\overrightarrow{PQ}=(x_2-x_1,\,y_2-y_1)$，
	因此上式就是数量积
	\[
	\underbrace{(A,B)}_{\text{系数向量}}\cdot
	\underbrace{(x_2-x_1,\,y_2-y_1)}_{\text{直线方向向量}}=0.
	\]
	两个非零向量数量积为零，说明它们垂直。因此 $(A,B)$ 垂直于直线 $l$，
	\textbf{它就是直线的一个法向量}：
	\[
	\boxed{Ax+By+C=0\quad\Longrightarrow\quad\vec n=(A,B).}
	\]
	这不是人为规定的公式，而是从\textbf{两点代入直线方程、再相减}直接推出来的。

	\medskip
	\textbf{\textcolor{CExplanation}{为什么需要法向量？（不用斜率也能判垂直）}}
	判断两直线垂直，也可以用斜率乘积为 $-1$，但竖直直线（如 $x=3$）
	\textbf{斜率不存在}，需要另外讨论。法向量直接由一般式的 $x,y$ 系数给出：
	$x=3$ 可写成 $1x+0y-3=0$，其法向量就是 $(1,0)$。
	所以法向量的作用是：\textbf{把几何上的垂直，转化成向量数量积为零；
	不需要求斜率，也不需要因竖直线而分类讨论。}
	对 $l_1,l_2$，只需检验 $(A_1,B_1)\cdot(A_2,B_2)=0$。

	\medskip
	\textbf{\textcolor{CExplanation}{几何本质（法向量垂直）}}
	垂直的几何含义是两条直线的方向相互成$90^\circ$角；通过法向量垂直来刻画最为简洁。

	\medskip
	\textbf{\textcolor{CExplanation}{代数意义（系数运算关系）}}
	垂直判断只需要计算两组系数的乘积之和，无需化为斜截式，适用于所有直线（包括斜率不存在的情况）。

	\medskip
	\textbf{\textcolor{CExplanation}{考点价值（快速判定垂直）}}
	\begin{itemize}[leftmargin=*, itemsep=0.5em, topsep=0.4em]
		\item \textbf{考法一（直接验证垂直）：}给出两条直线一般式方程，直接计算$A_1A_2+B_1B_2$的值，根据是否为零判断是否垂直。
		\item \textbf{考法二（求参数值）：}已知两条直线垂直，求方程中待定系数的值，利用$A_1A_2+B_1B_2=0$列方程求解。
	\end{itemize}

	\medskip
	\textbf{\textcolor{CExplanation}{顿悟点}}
	\textbf{关键：不要将垂直条件与平行或重合的条件（$A_1B_2-A_2B_1=0$）混淆；垂直条件中不涉及常数项$C$。}

	\medskip
	\textbf{\textcolor{CExplanation}{使用场景}}
	\begin{itemize}[leftmargin=*, itemsep=0.5em, topsep=0.4em]
		\item \textbf{场景一（已知方程判垂直）：}已知两条直线的具体方程，快速判断它们是否垂直。
		\item \textbf{场景二（求垂直直线方程）：}求过某点且与给定直线垂直的直线方程，可先设出一般式，利用垂直条件确定系数关系。
	\end{itemize}

\end{explanationbox}
